\subsection{A separability criterion}

PROPOSITION 11. --- Suppose that K is of characteristic exponent p, and let E be an algebraic extension of K, generated by a set S. If E is separable over K, then \(\mathrm{E} = \mathrm{K}(S^{p^{n}})\) for all integers \(n \geqslant 0\) ; conversely if E is of finite degree over K and \(\mathrm{E} = \mathrm{K}(S^{p})\) , then E is separable over K.

The case \(p = 1\) is trivial by the Cor. of V, p.~37. Suppose from now on that \(p \neq 1\) .

By hypothesis E is algebraic over K and we have \(E = K(S)\) , hence

\[
\mathrm{K} \left(\mathrm{S} ^ {p}\right) = \mathrm{K} \left(\mathrm{E} ^ {p}\right) = \mathrm{K} \left[ \mathrm{E} ^ {p} \right] \quad \text { by   V,   p.18,   Cor.   1 }.
\]

If E is of finite degree over K, it is a separable extension of K if and only if it is an etale algebra over K; the Cor. of V, p.~35 shows that this happens if and only if \(E = K[E^{p}]\) .

Suppose now that E is separable and of infinite degree over K. Then \(K[E^{p}]\) is the union of the subrings \(K[E^{\prime p}]\) where \(E'\) ranges over the set of subextensions of E of finite degree over K; but such an extension \(E'\) is separable over K (V, p.~36, Prop. 1), whence \(E' = K[E^{\prime p}] \subset K[E^{p}]\) by what has been said; finally we have \(E = K[E^{p}]\) . By induction on \(n \geqslant 0\) , the relation \(E = K[E^{p}]\) implies that \(E = K[E^{p^{n}}]\) .

COROLLARY 1. --- Every algebraic extension of a perfect field is a perfect field.

Let K be a perfect field of characteristic exponent p, and let E be an algebraic extension of K. Then E is separable over K (V, p.~36, Prop. 2) whence \(\mathrm{E} = \mathrm{K}(E^{p})\) by Prop. 11; but we have \(K = K^{p} \subset E^{p}\) , hence \(\mathrm{E} = \mathrm{K}(E^{p}) = E^{p}\) , and so E is perfect.

COROLLARY 2 (Mac Lane). --- Let \(\bar{\mathbf{K}}\) be an algebraic closure of \(\mathbf{K}\) and \(\mathrm{K}^{p^{-\infty}}\) the perfect closure of \(\mathbf{K}\) in \(\bar{\mathbf{K}}\) . For a subextension \(\mathbf{E}\) of \(\bar{\mathbf{K}}\) to be separable over \(\mathbf{K}\) it is necessary and sufficient that it should be linearly disjoint from \(\mathrm{K}^{p^{-\infty}}\) over \(\mathbf{K}\) .

We can immediately reduce to the case where \([E:K]\) is finite. Let \((x_{i})_{i\in I}\) be a basis of E over K. For E to be linearly disjoint from \(K^{p^{-\infty}}\) it is necessary and sufficient that it should be so from \(K^{p^{-n}}\) for all \(n\geqslant0\) , and this just means that the relation \(\sum_{i\in I}x_{i}a_{i}^{p^{-n}}=0\) implies \(a_{i}=0\) for all \(i\in I\) , for any family

\((a_{i})_{i\in I}\) of elements of K. This in turn means that the family \((x_{i}^{p^{n}})_{i\in I}\) is free over K, or also that it is a basis of the vector space E over K. Put differently, E is linearly disjoint from \(\mathbf{K}^{p^{-n}}\) if and only if \(\mathrm{E} = \mathrm{K}(\mathrm{E}^{p^n})\) . Now it suffices to apply Prop. 11.

Remarks. --- 1) When E is algebraic and of infinite degree over K, the condition \(E = K(E^{p})\) does not always ensure that E is separable over K. For example, if K is imperfect and E is a perfect closure of K, then we have \(E = K(E^{p})\) but E is not a separable extension of K (V, p.~39, Cor. 3).

\begin{enumerate}
\def\labelenumi{\arabic{enumi})}
\setcounter{enumi}{1}
\tightlist
\item
  Let \(\mathbf{E}\) be a separable algebraic extension of a field \(\mathbf{K}\) of characteristic exponent \(p\) . Then we have \(\mathbf{E}^p \cap \mathbf{K} = \mathbf{K}^p\) (Cor. 2); therefore if \(\mathbf{E}\) is perfect, the same is true of \(\mathbf{K}\) .
\end{enumerate}

